\begin{annexitem}{annex:p03.lq.1}{thm:p03-lq-representation}{Representation by class quantifiers}
\begin{proof}
Apply \(Q_K\) to the pattern's own atoms: \(R\bar x\) for each
relation, \(f(\bar x)=y\) for each function, and \(z=c\) for each
constant. The arguments always describe the given structure itself,
including valid function graphs and unique constants. The resulting
sentence is true exactly in \(K\), proving the first clause.

For the second, let \(\kappa\in\mathscr L[\sigma]\) define \(K\).
Induct on the construction of the extended formula. Base formulas
are already in \(\mathscr L\), and Boolean operations and
quantification use regularity. At a \(Q_K\)-application, first
translate all its operands by induction. Regard their free parameters
as fresh constants. Rename the finite pattern disjointly from the
ambient vocabulary, and substitute the translated operands for its
relations, function graphs and constant predicates in \(\kappa\).
On admissible interpretations the substitution equivalence gives
exactly the required truth value. Conjoin the total-single-valued
and singleton conditions for the finitely many graph and constant
operands. These conditions are formed from the translated operands
by quantification and Boolean operations. On inadmissible
interpretations their conjunction is false; on admissible ones it
is true. Thus the guarded substituted formula has the required
semantics on every structure, not merely on graph structures.
Fresh bound variables prevent capture, and the free parameters
remain unchanged. This completes the induction.

For the last clause, every formula of the first-order quantifier
extension is translated into \(\mathscr L\) by the same finite
induction, using the defining \(\mathscr L\)-sentence for each
quantifier actually occurring. Conversely, take a sentence
\(\varphi\in\mathscr L[\tau]\). Finitarity supplies a finite
\(\tau_0\subseteq\tau\) containing it. The class
\(K=\operatorname{Mod}^{\tau_0}_{\mathscr L}(\varphi)\)
is isomorphism-closed. The first clause defines it by its designated
quantifier, and the reduct property gives equivalence over \(\tau\).
For a formula with free variables, first use fresh constants for its
finite parameter list and include them in \(\tau_0\). In the
identity application just constructed, use the operand \(z=x\)
for the constant representing a free variable \(x\). This operand
always denotes the unique value assigned to \(x\). The same
semantic argument yields the formula equivalence under every
assignment, with no new free variables.

Full second-order logic satisfies the regularity hypotheses used here.
Its finite syntax gives finite occurrence; isomorphism, reduct and
renaming invariance follow by induction on formulas. Substitution is
capture-avoiding replacement on admissible graph interpretations.
For relativization to a nonempty closed domain, guard first-order
quantifiers by that domain, restrict quantified relations to tuples
from it, and require quantified functions to map its tuples into it.
Every function on the smaller domain extends to the larger universe
by a fixed value from the nonempty smaller domain; every smaller-domain
relation extends by false outside it. These extension facts prove both
directions of the second-order quantifier induction. Boolean operations
and first-order quantification are already present. Thus the representation
assertion applies to full second-order logic, as used below.
\end{proof}
\end{annexitem}

\begin{annexitem}{annex:p03.lq.2}{thm:p03-lq-second-order}{Second-order logic needs an infinite family}
\begin{proof}
Suppose, towards a contradiction, that the finite quantifier language
has the same expressive power as full second-order logic. The
identity application in \cref{annex:p03.lq.1} defines
each quantifier's class. Hence each such class has a fixed
second-order definition. Replace function and constant symbols in
each finite pattern by their graphs. Require the finitely many
graphs to be total and single-valued and the constant predicates
to be singletons. The resulting relational class is still
second-order definable, by
\cref{thm:p03-w03-p03-047-function-elimination}. Applying its
quantifier to the original graph operands has exactly the original
application's semantics, including rejection of malformed graphs.
The new graph-class operators are themselves second-order definable,
so their extension remains within second-order logic by
\cref{thm:p03-lq-representation}. It contains the original operator
language and hence would still be equivalent to second-order logic.
It therefore suffices to consider finite relational patterns with
fixed second-order defining sentences \(\kappa_1,\ldots,\kappa_s\).

We work in standard arithmetic \(\mathcal N\). We require truth
only for this fixed finite extension of first-order syntax, not for
all second-order formulas. We describe the local coding and truth
construction to make that distinction explicit. This coding changes
no arithmetic proof predicate or governing numbering elsewhere.

Use prefix strings over a finite alphabet for arithmetic operations,
equality, Boolean connectives, existential quantification, the
finitely many added quantifiers, delimiters, and variable tokens.
A variable index is written as a delimited unary natural number.
A quantifier node records the finite bound-variable lists and one
operand for each relation in its fixed pattern. A first-order
existential node records one bound variable and one operand.
These conventions determine an unambiguous finite grammar.
All abbreviations, including arithmetic graph predicates, are expanded
into that grammar before a code is taken.
Number the alphabet by positive integers. For a word
\(a_0\cdots a_{\ell-1}\), set
\[
 \operatorname{code}_Q(a_0\cdots a_{\ell-1})
   =2^\ell\prod_{j<\ell}p_{j+1}^{a_j+1},
\]
where \(p_0=2,p_1=3,\ldots\) are the successive primes.
The empty word has code \(1\).
Prime testing is bounded division; the next prime after \(p\)
can be found by bounded search up to \(p!+1\), since that number
has a prime divisor exceeding \(p\). Hence the prime sequence is
primitive recursive. Prime exponents are found by bounded search,
so word length and entries, concatenation, substring extraction,
and the test for a canonical word code are primitive recursive.
The code's length and every relevant search bound can be computed
primitive recursively from the input.

A left-to-right grammar check uses a finite stack of unfinished
nodes, their remaining operands, and their bound-variable lists.
At each symbol it either updates this coded stack or rejects.
It makes at most the word length many steps. This gives primitive
recursive term/formula recognition, child-operand extraction and
free-variable recognition. The same stack determines which
occurrences of a variable are free. Replacing each free occurrence
of \(x_0\) by the closed numeral for \(n\) requires no bound
variable renaming and produces a word whose length is bounded
by a primitive recursive function of the original length and \(n\).
Thus the total function \(d_Q\), which makes this substitution
on a formula code with free variables among \(x_0\) and uses the
fixed code of \(0=0\) on all other inputs, is primitive recursive.
These are the finite grammar operations underlying
\cref{thm:arithmetization-syntax}, now explicitly including the
finitely many new node forms; no semantics of the \(Q_i\) is
used by the algorithms.

Code finite sequences of natural numbers similarly, without the
finite-alphabet restriction, and read a sequence as an assignment
that is zero beyond its length. Reading a coordinate, extending
with zeros, and replacing finitely many coordinates are primitive
recursive. Term evaluation over \(\{0,S,+,\times\}\) is
primitive recursive: process the finitely many term nodes from
children to parents, storing the values already computed.
All these functions and their graph relations are first-order
definable in \(\mathcal N\), by
\cref{thm:p02-w02b-p02-034-pr-graph,thm:p05-rep-numeralwise-relations}.
Only their standard-model meanings are used.

Let \(T\) be a binary second-order relation on natural numbers,
intended to hold of a formula code and an assignment-sequence code.
There is a single second-order formula \(\operatorname{Sat}_Q(T)\)
with the following clauses, universally quantified over these codes.
Every compositional clause is guarded by well-formedness of the formula
and validity of the assignment code:
\begin{enumerate}
 \item On malformed formula or assignment codes, \(T\) is false.
       On an atomic formula, it has the truth value obtained by the
       primitive recursive term evaluations and arithmetic comparison.
 \item On a negation node it is the negation of the child's value;
       on a conjunction node it is the conjunction of the two child
       values; on an existential node binding \(x_j\), it holds
       exactly when the child holds at an assignment obtained by
       updating coordinate \(j\) to some natural number.
 \item At a node for \(Q_i\), interpret each relation of its pattern
       by the corresponding operand's \(T\)-truth set, with its
       bound coordinates updated to the proposed tuple of values.
       Substitute these truth-set formulas in the fixed second-order
       sentence \(\kappa_i\); require \(T\) at the node exactly
       when that substituted sentence is true.
\end{enumerate}
This is a finite second-order formula, not a schema: there are
finitely many node kinds and finitely many \(Q_i\). Node tests,
child codes, bound-coordinate lists and updates are expressed by
the first-order graph formulas just described. The third clause
uses one fixed finite formula for each \(i\); substitution for its
pattern relations is a finite syntactic operation, with bound
variables renamed to prevent capture of \(T\), the node code,
the assignment code and tuple parameters. Its second-order
quantifiers remain genuine second-order quantifiers over \(\mathcal N\).

The actual compositional truth relation for the fixed quantifier
language is a subset of \(\mathbb N^2\) and satisfies these clauses.
Existence follows by recursion on its finite formula trees. Conversely,
induction on a formula tree shows that any relation satisfying
\(\operatorname{Sat}_Q\) has exactly this value on every assignment:
the operands have already been determined before the quantifier
clause is evaluated. Malformed codes are fixed by the first clause.
Thus there is exactly one such relation in the full second-order
semantics. In particular the second-order formula
\[
 \operatorname{True}_Q(n)
   \ :=\ \operatorname{Sentence}_Q(n)\land
       \exists T\bigl(\operatorname{Sat}_Q(T)\land T(n,1)\bigr)
\]
defines truth of sentences of the fixed quantifier language in
\(\mathcal N\). The assignment code \(1\) is the empty sequence;
sentences have no free-variable values to supply.

The assumed expressive equivalence supplies a formula \(\psi(x_0)\)
in the finite quantifier language, over the same arithmetic vocabulary.
It is equivalent to \(\operatorname{True}_Q(x_0)\).
Let \(D_Q(x,z)\) be the first-order graph formula for \(d_Q\),
supplied by the primitive-recursive graph theorem. Form
\[
 \theta(x_0)=\exists z\bigl(D_Q(x_0,z)\land\neg\psi(z)\bigr),
 \qquad e=\operatorname{code}_Q(\theta),
\]
using fresh bound variables in the substitution for \(\psi\).
Let \(G=\theta(\bar e)\). By the definition of the local
substitution function, \(d_Q(e)=\operatorname{code}_Q(G)\).
The standard-model accuracy and uniqueness of its graph give
\[
 \mathcal N\models G
 \quad\Longleftrightarrow\quad
 \mathcal N\models\neg\psi\bigl(\overline{\operatorname{code}_Q(G)}\bigr).
\]
But \(G\) is a sentence of the very language whose truth \(\psi\)
defines. Its truth is therefore equivalent to the opposite right-hand
value, a contradiction. The proof neither assumes nor asserts a
second-order definition of second-order truth. It defines truth only
for a fixed finite quantifier language and uses the purported collapse
to produce the impossible self-definition.
\end{proof}
\end{annexitem}

\begin{annexitem}{annex:p03.lq.3}{thm:p03-lq-domain}{Domain translation and regularity}
\begin{proof}
The first assertion follows by taking the entire nonempty ambient
domain. We give both translations in the second assertion, including
the graph conditions.

Consider an application of \(Q_K^*\) with domain operand \(\delta\).
Write \(D(\bar x)=\bigwedge_j\delta(x_j)\); for the empty tuple
this is true. Construct a \(Q_{K^U}\)-application by assigning:
\begin{enumerate}
 \item to \(U\), the operand \(\delta(u)\);
 \item to an \(r\)-ary relation \(R\),
       \(D(\bar x)\land\alpha_R(\bar x)\);
 \item to an \(r\)-ary function \(f\), with \(r\geq1\), the graph
       \[
       \bigl(D(\bar x)\land\delta(y)\land\beta_f(\bar x,y)\bigr)
       \ \lor\
       \bigl(\neg D(\bar x)\land y=x_1\bigr);
       \]
 \item to a constant \(c\), \(\delta(z)\land\gamma_c(z)\).
\end{enumerate}
Treat a function of arity zero as a constant. Parameters, suppressed
in these displays, are retained throughout. Outside the selected
domain, the function graph has the explicit default value \(x_1\);
no element has to be chosen by the formula. Inside it, the graph
has exactly the values allowed by the original domain-restricted
operand. Therefore, if the original application is true, these
operands define total ambient functions, constants in \(D\), and
a nonempty closed \(D\) whose induced structure is precisely its
witness in \(K\). The resulting \(Q_{K^U}\)-application is true.
Conversely, if that application is true, its nonempty \(U\)-domain
is exactly \(D\), its functions are total and single-valued there,
and its constants lie there uniquely. The restricted operands are
exactly those of the original \(Q_K^*\)-application, which is
therefore true. This also proves that empty domains and malformed
graphs are rejected in both directions.

In the other direction, take an application of \(Q_{K^U}\), with
\(\delta\) its \(U\)-operand. Let \(G\) assert that each function
operand is a total single-valued graph on the whole ambient domain
and that each constant operand is a singleton there. This is a
finite conjunction of quantified formulas. Replace the application by
\[
 G\ \land\ Q_K^*\bigl(\delta;\alpha_R,\beta_f,\gamma_c\bigr),
\]
with the original bound lists understood. If the original application
holds, \(G\) holds and the closed subdomain supplies the relative
application. Conversely, under \(G\), each input tuple from \(D\)
has a unique ambient function value. Success of the relative
application gives a value in \(D\); ambient uniqueness forces it
to be that same value. Every constant likewise lies in \(D\).
Thus \(D\) is closed and the full-domain interpreted structure
belongs to \(K^U\). The factor \(G\) must not be omitted: a
relative application alone does not test function operands outside
its domain. Induction on the finite formula tree extends the two
application translations to whole formulas.

For regularity, isomorphism invariance, vocabulary extension, reducts
and renamings follow by induction from the class invariance and the
satisfaction definition. First-order atoms, Boolean operations and
quantifiers are part of the syntax. To relativize a formula to a
nonempty subdomain \(P\) closed under the ambient vocabulary, use
the usual guarded first-order quantifiers and recurse through the
Boolean operations. At a relative class quantifier, replace its
domain operand by
\[
 P(v)\land\delta^P(v),
\]
and recursively relativize every other operand. The tested domain
is now a subset of \(P\); satisfaction of every operand on its
tuples agrees, by induction, with satisfaction in the induced
structure on \(P\). Thus its interpreted pattern structure is
the same in both evaluations. The proof applies to a definable
domain as well: substitute its defining formula for \(P\), with
fresh bound variables. In particular the empty inner domain is
rejected on both sides.

Finally consider simultaneous substitution on admissible
interpretations. A formula has finite syntax and uses only finitely
many ambient symbols. Translate its terms by fresh variables naming
their values, using the prescribed graph and singleton formulas.
Use the value-formula recursion in
\cref{annex:p03.lq.8}, replacing each function graph
atom by its prescribed graph formula and each constant equality by
its prescribed singleton formula. Retained symbols use their own
graph atoms. On admissible interpretations the named values exist
uniquely; induction on terms proves that they are exactly the values
of the original terms. Translate atoms accordingly,
commute with Boolean operations and quantification, and translate
each class-quantifier operand recursively. For a relative
application, the translated domain and all the interpreted pattern
relations, functions and constants then agree with the originals
under the induced interpretation. The application truth values
agree by their definitions. This structural induction proves the
substitution property. Only finitely many templates occur in the
translated finite formula, even if the supplied replacement family
is infinite. Finite global graph guards can be added when needed
as in \cref{def:p03-w03-p03-047-relativization}.
\end{proof}
\end{annexitem}

\begin{annexitem}{annex:p03.lq.4}{thm:p03-lq-minimal}{The regular extension generated by a class}
\begin{proof}
Suppose a sentence \(\kappa\) defines \(K\) in a regular logic
\(\mathscr R\). Relativize \(\kappa\) to a fresh unary \(U\),
and conjoin the first-order conditions that \(U\) is nonempty,
contains each pattern constant and is closed under each pattern
function. There are finitely many such conditions. The resulting
sentence defines exactly \(K^U\), because its guard ensures the
hypothesis of the relativization equivalence. By
\cref{thm:p03-lq-representation}, every
\(\operatorname{FO}(Q_{K^U})\)-formula translates into
\(\mathscr R\). The domain translation of
\cref{thm:p03-lq-domain} therefore translates
\(\operatorname{FO}(Q_K^*)\) into \(\mathscr R\).
The latter logic is regular by that theorem and defines \(K\)
by the identity application on the whole domain, proving minimality.

For the more general assertion, translate base \(\mathscr L\)
formulas into a given regular extension and repeat the same
induction for each added application. This proves containment in
every regular extension defining \(K\). If
\(\mathscr L(Q_K^*)\) is itself regular, it is consequently the
least such extension. The hypothesis of regularity here is explicit;
finite Boolean closure of an arbitrary base logic is not by itself
a proof of all its substitution properties.
\end{proof}
\end{annexitem}

\begin{annexitem}{annex:p03.lq.5}{thm:p03-lq-chang-boundary}{A failure of relativization}
\begin{proof}
Construct two structures with \(V\subseteq U\). Their three disjoint
cells are \(V\), \(U\setminus V\), and the complement of \(U\).
In \(\mathfrak A\) give these cells cardinalities
\(\aleph_0,\aleph_1,\aleph_2\), respectively; in \(\mathfrak B\)
give them cardinalities \(\aleph_1,\aleph_1,\aleph_2\).
All cells are infinite. Each universe has cardinality \(\aleph_2\),
and the complement of \(U\) is its only cell of that cardinality.

We show that tuples taken from either of the two structures, including
twice from the same structure, satisfy the same formulas whenever their
equality patterns and cell memberships agree. The formulas are those of
\(\operatorname{FO}(Q^{\mathrm C})\). Induct on formulas. Atoms
record exactly that pattern; Boolean operations preserve agreement.
For an existential quantifier, a witness equal to an existing tuple
entry is matched by the corresponding entry. Any other witness can
be matched by a new element in its cell, since that cell is infinite
and only finitely many entries are excluded. Apply the induction
hypothesis to the enlarged tuple. The reverse transfer is identical.

For a Chang-quantifier step \(Q^{\mathrm C}x\,\psi(x,\bar a)\),
the induction hypothesis also applies to two assignments in the
same structure. It follows that \(\psi\) has constant truth value
on the non-parameter elements of each cell. Its truth set is
therefore a union of some cells with finitely many points removed,
together with some of the finitely many parameter values. This
set has cardinality \(\aleph_2\) exactly when the non-parameter
part of the complement of \(U\) is included. The induction
hypothesis gives the same inclusion on the other side. Hence the
Chang applications agree. This completes the induction, and the
empty tuple gives agreement on all sentences.

On the subdomain \(U\), however, \(Q^{\mathrm C}x\,V(x)\)
is false in \(\mathfrak A\), where \(|V|=\aleph_0<\aleph_1=|U|\),
and true in \(\mathfrak B\), where both cardinals are \(\aleph_1\).
There are no functions or constants to obstruct closure of \(U\).
If a sentence of \(\operatorname{FO}(Q^{\mathrm C})\) expressed
this relativization, it would distinguish the two structures,
contrary to the induction. Thus relativization fails.
\end{proof}
\end{annexitem}

\begin{annexitem}{annex:p03.lq.6}{thm:p03-lq-standard-omega}{Standard parts and projective representation}
\begin{proof}
For the first clause, conjoin the given first-order sentence with
an application of \(Q_K^*\) whose domain is \(U\) and whose
operands are the pattern's own relation atoms, function graphs
and constant equalities. The relative application requires exactly
that \(U\) be nonempty and closed and that its induced structure
belong to \(K\). This is precisely the permitted-structure
condition of the standard-part semantics.

For the second clause we give the projective construction. In a
permitted standard-order structure write \(N\) for \(U\), with
its fixed order \(<_N\) of type \(\omega\). The notation
\([0,n)_N\) abbreviates the definable set
\(\{j:U(j)\land j<_N n\}\). Let \(D(x,\bar a)\) and
\(R(x,y,\bar a)\) be any definable domain and binary relation.
All order conditions below concern only the restriction of \(R\)
to \(D\), not its values outside that domain.
The property that \((D,R)\) has order type \(\omega\) has
first-order projective witnesses in this standard-part semantics:
\begin{enumerate}
 \item A positive witness is a relation \(H^+(j,x,\bar a)\)
       giving an order isomorphism from \((N,<_N)\) onto \((D,R)\).
       Require it to be a bijection and require
       \[
       H^+(j,x,\bar a)\land H^+(k,y,\bar a)
          \ \longrightarrow\
       (j<_N k\leftrightarrow R(x,y,\bar a)).
       \]
 \item A negative witness is one of: failure of nonempty strict
       linear order on \(D\); a bijection from \([0,n)_N\)
       onto \(D\), for some \(n\in N\); or an injection from
       \(N\) into the predecessors of some \(d\in D\).
\end{enumerate}
These are finite first-order conditions on the displayed witness
relations and elements. To be explicit, a relation \(H\) is a
bijection between definable sets \(X,Y\) if it has no pairs outside
\(X\times Y\), every \(x\in X\) has exactly one \(H\)-successor
in \(Y\), and every \(y\in Y\) has exactly one \(H\)-predecessor
in \(X\). For an injection replace the last existence requirement
by at most one predecessor. All uniqueness conditions expand into
ordinary first-order quantifiers and equality.

The positive criterion is its own order-type definition. Each
negative criterion excludes order type \(\omega\). Conversely,
a nonempty strict linear order with no infinite predecessor set
is finite or has type \(\omega\), by the finite-predecessor
argument in \cref{annex:p03.general-pinning}. Hence a strict
linear order not of type \(\omega\) is finite or has an infinite
predecessor set. A finite set is bijective with exactly one
\([0,n)_N\), and every infinite set admits an injection from
\(N\) under Choice. Thus the negative criteria are exhaustive.
The empty domain is already included in failure of nonempty order.

Now take a fixed finite formula of
\(\operatorname{FO}(Q_{K_\omega}^*)\). Introduce, for each of
its subformulas, a relation recording its truth under its finite
list of free parameters. Impose the ordinary first-order
compositional clauses at atoms, Boolean connectives and
first-order quantifiers. At a relative quantifier node, use the
truth relations of its domain and order operands as \(D,R\).
When the node's truth relation holds, require a positive witness;
when it does not, require the disjunction of negative witness
conditions just given. Give the witness relations the node's
free parameter tuple as additional arguments. Introduce auxiliary
functions for the possible finite bound \(n\) and predecessor
point \(d\); impose their required membership only in the branch
where they are used. Values in unused branches are arbitrary.
There are only finitely many nodes and each parameter list is
finite, so the result is one first-order sentence over a finite
expansion of the formula's vocabulary. Require the truth predicate
for the original sentence to hold.

In every satisfying expansion, induction on subformulas proves
that the truth relations have their intended meanings: the two
quantifier witness tests are mutually exclusive and exhaustive.
Conversely, the actual truth relations satisfy the clauses after
choosing witnesses for all parameter tuples. Choice supplies these
set-indexed families, and their graphs are the auxiliary relations.
Default values for unused function arguments exist because the
ambient universe is nonempty. No second-order quantification has
been inserted into the sentence: the extra relations and functions
are exactly the symbols forgotten by the projective reduct.

If the target universe is infinite, choose a countably infinite
subset for \(N\) and order it with type \(\omega\). The
standard-part symbols are fresh, so this does not change the
target structure. The construction therefore gives precisely the
infinite target models as ordinary same-domain reducts of the
permitted expanded models.

For an arbitrary nonempty target universe \(A\), first enlarge
it by a disjoint countably infinite set carrying \(N\), and use
a fresh predicate \(V\) to mark \(A\). Extend the target
functions outside \(A\) by a fixed element of \(A\), keep its
constants in \(A\), and extend its relations arbitrarily outside
\(A\). Relativize the original sentence to \(V\), using
\cref{thm:p03-lq-domain}, then apply the preceding construction
on the enlarged, infinite universe. Require \(V\ne\varnothing\)
and closure under the finitely many function and constant symbols
actually occurring in the sentence. The relativized reduct to
\(V\) has exactly the original truth value.
If the target vocabulary contains further symbols, the RPC
semantics itself requires \(V\) to be closed under all its
functions and constants; in the forward construction the same
extension rule ensures this. Unused symbols require no additional
sentences. In the reverse direction they do not affect the
sentence, by the reduct property. This proves the all-cardinality
RPC assertion.

Finally every permitted \(\operatorname{FO}[K_\omega]\)-structure
contains an infinite subset \(N\), so its universe is infinite.
An ordinary reduct has the same universe. No such reduct can be
a one-element model of a tautology. This proves the third clause
and explains why the unrestricted all-cardinality assertion must
use relativized reducts, rather than silently treating PC as RPC.
\end{proof}
\end{annexitem}

\begin{annexitem}{annex:p03.lq.7}{thm:p03-lq-field-boundary}{Different compactness behaviours}
\begin{proof}
The permitted structures of \(\operatorname{FO}[K]\) are axiomatized
by an ordinary first-order theory: \(U\) is nonempty and closed
under the field vocabulary, its induced structure satisfies the
field axioms, and \(n\cdot1\ne0\) for every positive integer \(n\).
These are formulas with ordinary finite syntax; only the collection
of characteristic-zero axioms is infinite. If a set of ordinary
sentences is finitely satisfiable in permitted structures, then
its union with this fixed first-order theory is finitely
satisfiable. First-order compactness gives a model of the union,
which is again permitted and satisfies all the original sentences.
Thus the standard-part semantics is compact.

For the operator extension, use a fresh copy of the field
vocabulary and let \(\chi\) be the identity application of
\(Q_K^*\) to its field operations and constants with the whole
domain selected. On fields, \(\chi\) holds exactly in
characteristic zero. Consider the countable theory consisting of
the field axioms, \(\neg\chi\), and \(n\cdot1\ne0\) for
all positive \(n\). A finite part is satisfied in
\(\mathbb F_p(t)\) for a prime \(p\) larger than every integer
appearing in its characteristic requirements. This fraction field of the polynomial ring \(\mathbb F_p[t]\) is
infinite, since the powers of \(t\) are distinct, and has characteristic
\(p\). Thus \(\neg\chi\) also holds.
Primes of arbitrarily large size exist, for a prime divisor of
\(m!+1\) exceeds \(m\). The whole theory, however, would be a
field of characteristic zero satisfying \(\neg\chi\), which is
impossible. This proves the asserted failure of countable compactness.
\end{proof}
\end{annexitem}

\begin{annexitem}{annex:p03.lq.8}{thm:p03-lq-finite-comparison}{Finite comparison theorem}
\begin{proof}
We first give the term reduction and finiteness arguments needed for
both directions. A value formula for a term is defined recursively:
\[
\begin{aligned}
 \operatorname{Val}_x(y)&\ :=\ y=x,\\
 \operatorname{Val}_c(y)&\ :=\ c=y,\\
 \operatorname{Val}_{f(t_1,\ldots,t_r)}(y)
   &\ :=\ \exists z_1\cdots\exists z_r
       \left(\bigwedge_{j=1}^r\operatorname{Val}_{t_j}(z_j)
                    \land f(z_1,\ldots,z_r)=y\right).
\end{aligned}
\]
All introduced variables are fresh. Induction on terms shows that
\(\operatorname{Val}_t(y)\) holds exactly at the value of \(t\).
Replace an atom by existentially naming the values of its terms and
then applying its relation or equality to those variables. This
produces the indicated graph-atomic forms. Commute the translation
with Boolean operations and quantifiers, translating each operand
of a class quantifier separately. The interpreted relations in an
application are unchanged, so the translation preserves truth.
For an infinitary conjunction the same clause works componentwise,
with the original finite free-variable set retained. Quantifier
rank need not be preserved by this construction. When free variables
are renamed or identified in an infinitary formula, first injectively
rename its whole variable alphabet into the complement of the finite
list of intended free names. Then replace the renamed free variables
by those intended names. No intended free name is bound in the renamed
formula, so the replacements avoid capture even if the original formula
used every variable name somewhere. This renaming does not change rank.

For each finite number \(k\) of free variables and each \(m\),
there is a finite list of term-reduced formulas of rank at most
\(m\) from which every such formula is equivalent to one member.
Here equivalence is over all structures, not just the two under
comparison. At rank zero there are finitely many graph atoms in
\(k\) variables, because the vocabulary and all arities are finite.
Their Boolean combinations have finitely many truth-table forms.
Inductively, rank at most \(m+1\) formulas are Boolean combinations
of atoms, existential quantifications of rank at most \(m\)
formulas in \(k+1\) variables, and applications of the finitely
many class quantifiers to rank at most \(m\) operands. For an
operand binding \(r\) variables use the induction list in
\(k+r\) variables. Fresh bound-variable renaming reduces all
choices to these finite lists. Each pattern has finitely many
operands, so there are finitely many choices of their representatives.
Taking Boolean truth-table forms again proves the claim. No
decision procedure for semantic equivalence is assumed: the lists
may contain redundancies.

For a tuple \(\bar a\) in a structure \(\mathfrak A\), let
\(\vartheta^m_{\mathfrak A,\bar a}(\bar x)\) be the conjunction
of all members of such a finite list true at \(\bar a\) and the
negations of all members false there. Its rank is at most \(m\).
A tuple in any structure satisfies it exactly when it has the same
rank-at-most-\(m\) term-reduced type as \(\bar a\). In particular,
it fixes every equality and graph-atomic fact of the tuple.

\emph{From a system to formula agreement.}
Given an \(n\)-system, prove that a map \(p\in I_m\) preserves
all term-reduced formulas of rank at most \(m\) on its domain.
Induct on rank, using structural induction for Boolean combinations
at the same rank. Atoms are preserved by the definition of a graph
partial isomorphism. An existential witness is included by the
one-element extension property, after which the induction hypothesis
applies at level \(m-1\); back gives the reverse implication.

For a class-quantifier application true in \(\mathfrak A\), its
operand truth sets form a pattern structure \(\mathfrak C\in K\)
on \(A\). Use the quantifier-forth property to obtain
\(\mathfrak D\in K\) on \(B\). Every tuple of each relation of
\(\mathfrak D\) is matched, by a responding map at level \(m-1\),
to a tuple satisfying the corresponding operand in \(\mathfrak A\).
The induction hypothesis transfers that operand to \(\mathfrak B\).
Thus each relation of \(\mathfrak D\) is contained in the appropriate
operand truth set on \(B\). Monotonicity, on the same domain \(B\),
implies that these truth sets form a member of \(K\); hence the
application is true in \(\mathfrak B\). Use back for the converse.
Choose any map in the nonempty \(I_n\) and apply this assertion
to sentences, which have no assigned free variables. This proves
(2) implies (1).

\emph{From formula agreement to a system.}
For \(m\leq n\), let \(I_m\) consist of finite injective maps
whose domain enumeration \(\bar a\) and image enumeration
\(\bar b\) have the same rank-at-most-\(m\) type. This does not
depend on the enumeration: bijective renaming of free variables
preserves the reduced syntax and rank. The rank-zero atoms ensure
that these maps are graph partial isomorphisms. Sentence agreement
through rank \(n\) puts the empty map in every \(I_m\).

Take \(p\in I_{m+1}\), with enumerations \(\bar a,\bar b\),
and \(a\in A\). The formula
\(\exists y\,\vartheta^m_{\mathfrak A,\bar a a}(\bar x,y)\)
is true at \(\bar a\) and has rank at most \(m+1\). It is
therefore true at \(\bar b\); choose a witness \(b\).
The equality conjuncts ensure that adding \(a\mapsto b\) gives
an injective function extending \(p\). The characteristic formula
ensures that this extension lies in \(I_m\). If repeated coordinates
are deleted, any formula on the resulting domain enumeration can
be renamed back to the longer tuple without increasing rank, so
the definition of \(I_m\) is still met. This proves forth;
interchange the structures for back.

Now take \(\mathfrak C\in K\) on \(A\). For each pattern relation
\(R\) of arity \(r\), form
\[
 \theta_R(\bar x,\bar y)
   =\bigvee_{\bar c\in R^{\mathfrak C}}
          \vartheta^m_{\mathfrak A,\bar a\bar c}(\bar x,\bar y).
\]
There are only finitely many characteristic formulas in the fixed
number of variables, so this is a finite disjunction after removing
repetitions. An empty disjunction is false. On \(\mathfrak A\)
with parameters \(\bar a\), its truth set contains
\(R^{\mathfrak C}\). Monotonicity therefore makes the
\(Q_K\)-application to the \(\theta_R\) true at \(\bar a\).
Its rank is at most \(m+1\), so it is true at \(\bar b\).
Let \(\mathfrak D\) be the operand truth sets there; then
\(\mathfrak D\in K\). Every tuple \(\bar d\in R^{\mathfrak D}\)
satisfies the characteristic formula for some
\(\bar c\in R^{\mathfrak C}\). Equalities and type agreement give
an extension \(q\in I_m\) with \(q(\bar c)=\bar d\), exactly
as in the one-element case. This proves quantifier-forth, and
interchanging the structures proves back. The construction also
handles a nullary relation: its tuple is empty and the responding
map adds no entries. Hence \((I_0,\ldots,I_n)\) is an \(n\)-system.

Finally, full sentence equivalence puts the empty map in these same
type-defined \(I_m\) for every \(m<\omega\). The argument just
given works at every level, producing one \(\omega\)-system, not
merely unrelated finite systems. Conversely its finite initial
segments are \(n\)-systems. Systems at every finite level imply
agreement on every reduced sentence, because each has some finite
rank, and term reduction then gives agreement on all sentences.
This proves all the assertions.
\end{proof}
\end{annexitem}

\begin{annexitem}{annex:p03.lq.14}{thm:p03-lq-characteristic}{Characteristic formulas and types}
\begin{proof}
We induct on \(m\). At zero the finite graph-atomic diagram
specifies exactly the rank-zero type, including all Boolean
combinations of graph atoms, and has rank zero.
Assume the assertion for rank \(m\). The finite-basis argument
in \cref{annex:p03.lq.8} applies for each fixed number of free
variables. Hence only finitely many equivalence classes of the
\(\xi^m\)-formulas occur in each displayed disjunction. There
are then only finitely many different disjunctions up to equivalence,
and only finitely many corresponding quantified conjuncts.
Choose an actual operand representing each occurring class.
All actual \(\xi^m\) operands have rank \(m\) by induction.
Since \(A\ne\varnothing\), the existential and universal
one-element clauses retain a quantified operand of rank \(m\).
The displayed formula therefore has rank exactly \(m+1\).

It is true at \(\bar a\). The one-element clauses use each
\(c\) as its own witness. In a positive quantifier clause,
the truth set of \(O^m_{\bar a,M}\) contains \(M\), so
monotonicity gives membership in \(K\) when \((A,M)\in K\).
For a negative clause, the truth set of
\(\neg O^m_{\bar a,A^r\setminus M}\) is contained in \(M\):
every tuple outside \(M\) satisfies its own characteristic
formula in that disjunction. If this smaller truth set belonged
to \(K\), monotonicity would put \(M\) in \(K\), contrary to
the index of the conjunct. Thus the negative clause is also true.

Suppose \(\mathfrak B\models\xi^{m+1}_{\bar a}(\bar b)\).
The universal one-element clause, and nonemptiness of \(B\),
show that the base tuples have the same graph-atomic type: choose
one element on the \(B\)-side and its matching \(\xi^m\)-type
from the disjunction, then restrict to the base variables.
An existential formula with matrix of rank at most \(m\)
transfers from \(\mathfrak A\) by the corresponding existential
characteristic conjunct. It transfers back by the universal
one-element conjunct and the induction hypothesis. Boolean
operations preserve agreement.

For a quantifier application \(Q_K\bar y\,\rho(\bar x,\bar y)\)
with matrix of rank at most \(m\), let \(M\) be the matrix's
truth set in \(\mathfrak A\) at \(\bar a\).
If \((A,M)\in K\), the positive conjunct supplies
\(Q_K\bar y\,O^m_{\bar a,M}\) in \(\mathfrak B\).
Every tuple in that disjunction has the type of some tuple of
\(M\), so satisfies \(\rho\) by induction. Monotonicity
therefore gives the required positive application in \(\mathfrak B\).
If \((A,M)\notin K\), the negative conjunct says that the
truth set of \(\neg O^m_{\bar a,A^r\setminus M}\) in
\(\mathfrak B\) does not belong to \(K\).
The truth set of \(\rho\) is contained in this set: a tuple
matching the type of one outside \(M\) must make \(\rho\)
false. Hence monotonicity prevents the truth set of \(\rho\)
from belonging to \(K\). This gives the negative application.
Structural induction now proves agreement on every formula of
rank at most \(m+1\).

Conversely, any tuple having that rank type satisfies
\(\xi^{m+1}_{\bar a}\), since this is a formula of the same
rank true at \(\bar a\). This proves the characteristic-type
claim. Restricting a tuple type to fewer free variables proves
the projection entailment, and restricting a rank-\((m+1)\)
type to rank at most \(m\) proves the descending entailment.
\end{proof}
\end{annexitem}

\begin{annexitem}{annex:p03.lq.9}{thm:p03-lq-karp}{Infinitary comparison theorem}
\begin{proof}
Suppose first that an unbounded system \(I\) is given. After term
reduction, structural induction on well-founded formula trees proves
preservation by every \(p\in I\). Atoms use graph partial isomorphism;
negation and every set conjunction use their truth clauses.
Existential witnesses use a responding extension in the same \(I\).
At a class quantifier, use the quantifier-forth condition and
monotonicity exactly as in
\cref{annex:p03.lq.8}; the induction hypothesis
applies to the smaller operand formulas at the responding maps.
Back proves the reverse directions. Any member of the nonempty
system therefore gives agreement on sentences.

Conversely, suppose the structures agree on all sentences of the
infinitary language. Let \(I\) consist of finite maps for which a
domain enumeration and its image satisfy exactly the same formulas
of that language. Equality atoms make these maps injective and
well-defined, and the other graph atoms make them graph partial
isomorphisms. The empty map belongs to \(I\).
We construct the required extensions using set-indexed formulas;
we do not conjoin the proper class of all formulas.

Fix \(p\in I\), with enumerations \(\bar a,\bar b\), and a
finite tuple \(\bar c\) from \(A\). Call a tuple \(\bar d\)
of the same length from \(B\) bad if
\((\bar a,\bar c)\) and \((\bar b,\bar d)\) fail to agree
on all formulas. For each bad tuple choose a distinguishing formula
\(\rho_{\bar d}(\bar x,\bar y)\), negated if necessary, true
at \((\bar a,\bar c)\) and false at \((\bar b,\bar d)\).
Each witness has free variables among the displayed finite
list. Put
\[
 \zeta_{\bar c}(\bar x,\bar y)
       =\bigwedge_{\bar d\text{ bad}}\rho_{\bar d}(\bar x,\bar y).
\]
This is a legitimate set conjunction, indexed by a subset of a
finite power of \(B\). It is true at \((\bar a,\bar c)\).
If it holds at \((\bar b,\bar d)\), that tuple is not bad, so
adjoining \(\bar c\mapsto\bar d\) gives a member of \(I\)
extending \(p\). Equalities in the common type guarantee
well-definedness and injectivity even when tuple entries repeat.
Deleting repeated coordinates does not lose type agreement, by
renaming free variables. Empty conjunctions are true.
The choices of formulas are set-sized: for each index one can
first select the least possible set-theoretic rank of a witness,
then apply Choice within the resulting sets of witnesses.

For a single \(c\in A\), the formula
\(\exists y\,\zeta_c(\bar x,y)\) is true at \(\bar a\),
so it is true at \(\bar b\). A witness gives an extension in
\(I\) containing \(c\). The symmetric construction gives back.

For quantifier-forth, let \(\mathfrak C\in K\) have domain \(A\).
For each pattern relation \(R\), put
\[
 \eta_R(\bar x,\bar y)
      =\bigvee_{\bar c\in R^{\mathfrak C}}
                    \zeta_{\bar c}(\bar x,\bar y).
\]
This is again a set-indexed formula with a finite free-variable
list. In \(\mathfrak A\), at \(\bar a\), its truth set contains
\(R^{\mathfrak C}\). Monotonicity makes the application of
\(Q_K\) to these operands true at \(\bar a\).
Agreement of the base tuples on all formulas transfers it to
\(\bar b\). Its operand truth sets on \(B\) therefore form
\(\mathfrak D\in K\). A tuple in \(R^{\mathfrak D}\) satisfies
some \(\zeta_{\bar c}\) with \(\bar c\in R^{\mathfrak C}\),
and hence has the required responding extension in \(I\).
The same construction with the structures interchanged gives
quantifier-back. It works for each quantifier separately, so
neither the ambient vocabulary nor the quantifier family has to
be finite. This proves that \(I\) is an unbounded system.
\end{proof}
\end{annexitem}

\begin{annexitem}{annex:p03.lq.10}{thm:p03-lq-countable}{Ordinary comparison systems}
\begin{proof}
The preservation induction in \cref{annex:p03.lq.9}, together
with term reduction, shows that every map in an unbounded graph
system preserves every atomic formula with arbitrary terms.
Its forth and back conditions are therefore exactly those of an
ordinary partial isomorphism system. Apply
\cref{thm:p03-w03-p03-039-countable-back-forth} for the countable
conclusion. Explicitly, alternate extending the map to include
successive elements of enumerations of the two structures. The
union of the increasing maps is an everywhere-defined bijection;
each finite atomic instance is already decided at some stage,
so the union is an isomorphism.
If two structures agree in \(L_{\infty\omega}\),
\cref{thm:p03-w03-p03-039-karp} supplies an ordinary partial
isomorphism system. A logic \(\mathscr R\) with the stated
preservation property then agrees on them. This is exactly
\(\mathscr R\leq_{\equiv}L_{\infty\omega}\).

The finite increasing partial bijections between \((\mathbb Q,<)\)
and \((\mathbb R,<)\) form an unbounded system: a new point lies
in one of the finitely many intervals cut out by the current domain,
and the corresponding interval on the other side is nonempty by
density and absence of endpoints. A point not already used can be
chosen there. The same argument works in reverse. The two orders
have different cardinalities and hence are not isomorphic. Thus
the countability hypothesis in the isomorphism conclusion is essential.
\end{proof}
\end{annexitem}

\begin{annexitem}{annex:p03.lq.11}{thm:p03-lq-qone-game}{Games and uncountability}
\begin{proof}
First compare an \(n\)-system with the \(n\)-round game. To make
the initial empty position available, replace each \(I_m\) by all
finite restrictions of its members. These sets remain nonempty,
contain the empty map, and have the extension properties: given a
restriction \(p\) of an old map \(p'\), apply the old property to
\(p'\) and restrict its answering extension to the domain of
\(p\) together with the finitely many chosen response elements.
The range version is obtained by including the corresponding
preimages. The resulting restriction belongs to the new lower level.

The second player maintains a map at the level equal to the number
of rounds remaining, extending all pairs actually selected. An
ordinary move uses forth or back. For an uncountability move,
quantifier-forth or back supplies an uncountable answering set;
every point then selected by the first player has a mate and a
responding extension at the lower level. At level zero, the map is
a graph partial isomorphism, and so is its restriction to the played
pairs. Choice permits fixing the responses as a strategy.

Conversely, let \(I_m\) consist of all finite graph partial
isomorphisms from which the second player can win the remaining
\(m\) rounds. A strategy winning from the empty map for \(n\)
rounds supplies nonempty \(I_m\) at every \(m\leq n\): follow
any legal ordinary moves for the preceding rounds and take the
resulting position. The current graph condition cannot have been
violated at a winning position. A winning response to an ordinary
move has a winning continuation for \(m\) rounds, giving forth
and back from \(I_{m+1}\) to \(I_m\). A winning response to an
uncountability move supplies the required uncountable answering
set. For each point that the first player may then select, the
strategy supplies a mate with a winning \(m\)-round continuation.
These are exactly the quantifier extension conditions. Thus a
winning strategy for \(n\) rounds gives an \(n\)-system.
Apply \cref{thm:p03-lq-finite-comparison} at all finite levels.
The term-reduction clause there accounts for arbitrary terms in
sentences, without asserting that it preserves a fixed rank.
\end{proof}
\end{annexitem}

\begin{annexitem}{annex:p03.lq.12}{thm:p03-lq-qone-interpolation}{Failure of interpolation for uncountability}
\begin{proof}
Let \(I\) be all finite injective maps preserving and reflecting
the equivalence relation. The empty map belongs to \(I\).
For an ordinary new element in an already represented class,
choose a fresh element in its matched class. For an element in a
new class, choose a fresh element of a class not yet represented
on the other side. Each class is infinite and there are infinitely
many classes, so both choices are possible. This proves ordinary
forth and back.

For quantifier-forth fix \(p\in I\) and an uncountable
\(C\subseteq A\). Partition \(C\) into its intersections with
the finitely many classes represented in the domain of \(p\),
and its part outside all those classes. At least one part is
uncountable, since a finite union of countable sets is countable.
If an intersection with a represented class is uncountable, take
\(D\) to be the matched class in \(B\), omitting the finite
range of \(p\). If the outside part is uncountable, take \(D\)
to be any class in \(B\) not yet represented. In either case
\(D\) is uncountable. Each \(b\in D\) can be matched with a
fresh \(a\) from the chosen uncountable part of \(C\); adding
\(a\mapsto b\) extends \(p\) to a member of \(I\).
This is the \(Q_1\)-forth condition. The identical partition
argument on \(B\) gives back. Thus \(I\) is an unbounded system,
and \cref{thm:p03-lq-karp} gives the claimed infinitary equivalence.

For the interpolation counterexample, take
\(A=\omega\times\omega_1\) and \(B=\omega_1\times\omega_1\),
with equivalence given by equality of the first coordinate.
All classes have cardinality \(\aleph_1\); there are countably
infinitely many in \(\mathfrak A\) and uncountably many in
\(\mathfrak B\). Expand \(\mathfrak A\) by the countable
representative set \(R=\omega\times\{0\}\), and expand
\(\mathfrak B\) by the uncountable representative set
\(S=\omega_1\times\{0\}\). These expansions satisfy
\(\varphi_0(E,R)\) and \(\varphi_1(E,S)\), respectively, in
\cref{def:p03-w03-p03-047-keisler-entailment}.

Suppose an interpolant \(\theta(E)\) existed for
\(\varphi_0\models\neg\varphi_1\).
The first expansion and \(\varphi_0\models\theta\) give
\(\mathfrak A\models\theta\). The equivalence just proved
then gives \(\mathfrak B\models\theta\). Its second expansion
would satisfy both \(\theta\) and \(\varphi_1\), contrary to
\(\theta\models\neg\varphi_1\). This proves failure of
interpolation. In fact the same argument excludes an interpolant
in \(L_{\infty\omega}(Q_1)\).
\end{proof}
\end{annexitem}

\begin{annexitem}{annex:p03.lq.13}{thm:p03-lq-examples}{Examples of ordinary comparison systems}
\begin{proof}
For infinite sets, take all finite bijections between subsets. A
fresh element on either side can be matched to a fresh element on
the other, so forth and back hold. Equality is the only atomic
condition.

For dense linear orders without endpoints, take all finite
increasing partial bijections. The finitely many selected points
cut each order into corresponding open intervals and end intervals.
Density and absence of endpoints supply a new point in every
required interval. This gives forth and back and preserves all
order atoms.

To obtain nonisomorphic examples of any uncountable cardinality
\(\kappa\), let \(L\) be \(\kappa\) copies of \(\mathbb Q\)
ordered lexicographically by the ordinal \(\kappa\), and let
\(L'\) use the reverse ordering of the indices, leaving each copy
of \(\mathbb Q\) in its usual order. Both are dense without
endpoints and have cardinality \(\kappa\).
The initial segment of \(L\) formed by indices below \(\omega_1\)
has no countable cofinal subset: the indices of a countable subset
are bounded below \(\omega_1\). It may be the whole order when
\(\kappa=\aleph_1\), which causes no difficulty.
In contrast, every nonempty initial segment \(J\) of \(L'\) has
a countable cofinal subset. The set of indices of blocks meeting
\(J\) has a least ordinal \(\alpha\). The nonempty part of the
\(\alpha\)-block in \(J\) is countable and is cofinal in \(J\),
because all larger indices precede that block in \(L'\).
An isomorphism preserves initial segments and the existence of a
countable cofinal subset. Hence \(L\) and \(L'\) are not isomorphic.
Their finite increasing partial bijections nevertheless form the
system already constructed.

For the fields, identify their isomorphic prime fields. Let \(I\)
be all finite restrictions of field isomorphisms between subfields
finitely generated over those prime fields. It contains the empty
map. Every member preserves atomic formulas with arbitrary field
terms, since its witnessing field isomorphism preserves polynomial
and rational operations.

Let \(p\in I\) be witnessed by \(h:F\to G\), and take an element
\(a\) of the first field. If \(a\) is algebraic over \(F\),
let its monic minimal polynomial be \(r\). The polynomial obtained
by applying \(h\) to its coefficients is irreducible over \(G\)
and has a root \(b\) in the second algebraically closed field.
The evaluation maps identify \(F[a]\) and \(G[b]\) with the
respective polynomial quotients by \(r\) and its image. Thus
\(h\) extends to an isomorphism \(F(a)\to G(b)\) sending
\(a\) to \(b\).
If \(a\) is transcendental over \(F\), choose \(b\)
transcendental over \(G\). Such a \(b\) exists because \(G\)
has finite transcendence degree and the second field has infinite
transcendence degree over its prime field. Sending \(a\) to
\(b\) extends \(h\) to the rational function fields \(F(a)\)
and \(G(b)\): polynomial relations are preserved because neither
new element satisfies a nonzero polynomial over its base field.
In either case these enlarged subfields are still finitely
generated. Restricting the extension to the old finite domain
and \(a\) gives a member of \(I\) extending \(p\).
Apply the same argument to \(h^{-1}\) for back. This proves
all three assertions.
\end{proof}
\end{annexitem}

